% problem_id: S001-theorem-DM
% source_id: S001 (Stacks Project)
% source_file: stacks-morphisms.tex
% statement_locator: stacks-morphisms.tex lines 4375--4384
% proof_locator: stacks-morphisms.tex lines 4386--4526
% env: theorem
% id_note: label 在本来源内唯一
% status: complete (statement + author proof verbatim + reason + locators)
%
\begin{theorem}
\label{theorem-DM}
Let $\mathcal{X}$ be an algebraic stack. The following are equivalent
\begin{enumerate}
\item $\mathcal{X}$ is DM,
\item $\mathcal{X}$ is Deligne-Mumford, and
\item there exists a scheme $W$ and a surjective \'etale
morphism $W \to \mathcal{X}$.
\end{enumerate}
\end{theorem}

\begin{proof}
Recall that (3) is the definition of (2), see
Algebraic Stacks, Definition \ref{algebraic-definition-deligne-mumford}.
The implication (3) $\Rightarrow$ (1) is
Lemma \ref{lemma-properties-covering-imply-diagonal}.
Assume (1). Let $x \in |\mathcal{X}|$ be a finite type point.
We will produce a scheme over $\mathcal{X}$ which ``works'' in a
neighbourhood of $x$. At the end of the proof we will take the disjoint
union of all of these to conclude.

\medskip\noindent
By
Lemma \ref{lemma-point-finite-type-monomorphism}
the residual gerbe $\mathcal{Z}_x$ of $\mathcal{X}$ at $x$ exists and
$\mathcal{Z}_x \to \mathcal{X}$ is locally of finite type. By
Lemma \ref{lemma-separation-properties-residual-gerbe}
the algebraic stack $\mathcal{Z}_x$ is DM. By
Lemma \ref{lemma-DM-residual-gerbe}
there exists a field $k$ and a surjective \'etale morphism
$z : \Spec(k) \to \mathcal{Z}_x$.
In particular the composition $x : \Spec(k) \to \mathcal{X}$
is locally of finite type (by
Morphisms of Spaces, Lemmas
\ref{spaces-morphisms-lemma-composition-finite-type} and
\ref{spaces-morphisms-lemma-etale-locally-finite-type}).

\medskip\noindent
Pick a scheme $U$ and a smooth morphism $U \to \mathcal{X}$ such that
$x$ is in the image of $|U| \to |\mathcal{X}|$.
Consider the following fibre square
$$
\xymatrix{
U \ar[d] & F  \ar[l] \ar[d] \\
\mathcal{X} & \Spec(k) \ar[l]_-x
}
$$
in other words $F = U \times_{\mathcal{X}, x} \Spec(k)$. By
Properties of Stacks, Lemma \ref{stacks-properties-lemma-points-cartesian}
we see that $F$ is nonempty.
As $\mathcal{Z}_x \to \mathcal{X}$ is a monomorphism we have
$$
\Spec(k) \times_{z, \mathcal{Z}_x, z} \Spec(k)
=
\Spec(k) \times_{x, \mathcal{X}, x} \Spec(k)
$$
with \'etale projection maps to $\Spec(k)$ by construction of $z$.
Since
$$
F \times_U F =
(\Spec(k) \times_\mathcal{X} \Spec(k))
\times_{\Spec(k)} F
$$
we see that the projections maps $F \times_U F \to F$ are \'etale as well.
It follows that $\Delta_{F/U} : F \to F \times_U F$ is \'etale (see
Morphisms of Spaces, Lemma \ref{spaces-morphisms-lemma-etale-permanence}).
By
Morphisms of Spaces, Lemma
\ref{spaces-morphisms-lemma-etale-universally-injective-open}
this implies that $\Delta_{F/U}$ is an open immersion, which finally
implies by
Morphisms of Spaces, Lemma
\ref{spaces-morphisms-lemma-diagonal-unramified-morphism}
that $F \to U$ is unramified.

\medskip\noindent
Pick a nonempty affine scheme $V$ and an \'etale morphism $V \to F$.
(This could be avoided by working directly with $F$, but it seems easier
to explain what's going on by doing so.) Picture
$$
\xymatrix{
U \ar[d] & F  \ar[l] \ar[d] & V \ar[l] \ar[ld] \\
\mathcal{X} & \Spec(k) \ar[l]_-x
}
$$
Then $V \to \Spec(k)$ is a smooth morphism of schemes and $V \to U$ is an
unramified morphism of schemes (see
Morphisms of Spaces, Lemmas
\ref{spaces-morphisms-lemma-composition-smooth} and
\ref{spaces-morphisms-lemma-composition-unramified}).
Pick a closed point $v \in V$ with $k \subset \kappa(v)$ finite separable, see
Varieties, Lemma \ref{varieties-lemma-smooth-separable-closed-points-dense}.
Let $u \in U$ be the image point. The local ring
$\mathcal{O}_{V, v}$ is regular (see
Varieties, Lemma \ref{varieties-lemma-smooth-regular})
and the local ring homomorphism
$$
\varphi : \mathcal{O}_{U, u} \longrightarrow \mathcal{O}_{V, v}
$$
coming from the morphism $V \to U$ is such that
$\varphi(\mathfrak m_u)\mathcal{O}_{V, v} = \mathfrak m_v$, see
Morphisms, Lemma \ref{morphisms-lemma-unramified-at-point}.
Hence we can find $f_1, \ldots, f_d \in \mathcal{O}_{U, u}$
such that the images $\varphi(f_1), \ldots, \varphi(f_d)$
form a basis for $\mathfrak m_v/\mathfrak m_v^2$ over $\kappa(v)$.
Since $\mathcal{O}_{V, v}$ is a regular local ring this implies
that $\varphi(f_1), \ldots, \varphi(f_d)$ form a regular sequence
in $\mathcal{O}_{V, v}$ (see
Algebra, Lemma \ref{algebra-lemma-regular-ring-CM}).
After replacing $U$ by an open neighbourhood of $u$ we may assume
$f_1, \ldots, f_d \in \Gamma(U, \mathcal{O}_U)$. After replacing
$U$ by a possibly even smaller open neighbourhood of $u$ we may
assume that $V(f_1, \ldots, f_d) \to \mathcal{X}$ is flat and
locally of finite presentation, see
Lemma \ref{lemma-slice}.
By construction
$$
V(f_1, \ldots, f_d) \times_\mathcal{X} \Spec(k)
\longleftarrow
V(f_1, \ldots, f_d) \times_U V
$$
is \'etale and $V(f_1, \ldots, f_d) \times_U V$
is the closed subscheme $T \subset V$ cut out by $f_1|_V, \ldots, f_d|_V$.
Hence by construction $v \in T$ and
$$
\mathcal{O}_{T, v} =
\mathcal{O}_{V, v}/(\varphi(f_1), \ldots, \varphi(f_d)) = \kappa(v)
$$
a finite separable extension of $k$. It follows that $T \to \Spec(k)$
is unramified at $v$, see
Morphisms, Lemma \ref{morphisms-lemma-unramified-at-point}.
By definition of an unramified morphism of algebraic spaces this means that
$V(f_1, \ldots, f_d) \times_\mathcal{X} \Spec(k) \to \Spec(k)$
is unramified at the image of $v$ in
$V(f_1, \ldots, f_d) \times_\mathcal{X} \Spec(k)$.
Applying
Lemma \ref{lemma-etale-at-point}
we see that on shrinking $U$ to yet another open neighbourhood of $u$
the morphism $V(f_1, \ldots, f_d) \to \mathcal{X}$ is \'etale.

\medskip\noindent
We conclude that for every finite type point $x$ of $\mathcal{X}$ there
exists an \'etale morphism $f_x : W_x \to \mathcal{X}$ with $x$ in the
image of $|f_x|$. Set $W = \coprod_x W_x$ and $f = \coprod f_x$. Then $f$
is \'etale. In particular the image of $|f|$ is open, see
Properties of Stacks, Lemma \ref{stacks-properties-lemma-topology-points}.
By construction the image contains all finite type points of $\mathcal{X}$,
hence $f$ is surjective by
Lemma \ref{lemma-enough-finite-type-points} (and
Properties of Stacks, Lemma
\ref{stacks-properties-lemma-characterize-surjective}).
\end{proof}
